\begin{helper}[Early repeated-color normalization]
Let \((\varphi,T)\) be a minimum bad pair with
\(2<q:=\rho(T)<n:=|E(T)|\), and suppose that no color missing at
\(y_n\) is missing at any \(y_s\) with \(q\le s<n\).  Then there are a
coloring \(\varphi'\), a color \(\beta\), and an index \(i\) with
\(i+2\le q\) such that \((\varphi',T)\) is a minimum bad pair,
\(\beta\) is missing at both \(y_i\) and \(y_n\) under \(\varphi'\),
and the same no-repetition condition holds under \(\varphi'\).
\end{helper}
\begin{proof}
Unfolding \noderef{TashTreeMinimumBadPair} and
\noderef{TashTreeBadPair}, the full vertex set of \(T\) is not
elementary, and the bad-pair conjunct contains a \(k\)-triple.
By \noderef{KTriple}, \(\varphi\) is proper; this is the properness used
in every incidence-component argument below.  The exact penultimate
prefix supplied by \noderef{TreeSequencePrefix} is elementary by
\noderef{TashTreeProperPrefixElementary}.  Non-elementariness of the
full set, expanded with \noderef{ColoredVertexSetElementary}, gives two
distinct displayed vertices sharing a missing color.  By
\noderef{TreeSequenceVertexSet} every displayed vertex except \(y_n\)
belongs to the penultimate prefix, and \noderef{TreeSequence} makes
vertices at distinct displayed indices distinct.  Hence the two
witnesses cannot both be in that elementary prefix: one is \(y_n\) and
the other is some \(y_i\), \(i<n\).  Denote their common missing color
by \(\beta\).  The assumed no-repetition condition forces \(i<q\).
If \(i+2\le q\), take \(\varphi'=\varphi\); all the asserted
conclusions are immediate.

It remains to treat \(i+2>q\).  Since \(i<q\), this means
\(i=q-1\).  As \(q>2\), the prefix \(T_{q-2}\) has positive length and
is proper.  By \noderef{TashTreeFourUnusedColors}, it has at least four
colors in \noderef{TreeFreeColors}.  Each such color is missing at a
displayed vertex of \(T_{q-2}\), by
\noderef{TreeMissingColors}, and occurs on no colored exact edge of that
prefix, by \noderef{TreeUsedColors}.  Passing from \(T_{q-2}\) to
\(T_{q-1}\) adds only the exact edge \(e_{q-1}\), so at most one of the
four colors ceases to be unused.  Excluding also \(\beta\), choose a
remaining color \(\alpha\ne\beta\), missing at some \(y_m\) with
\(m\le q-2\), which is unused through \(y_{q-1}\).

The color \(\beta\) is unused through \(y_{q-1}\) as well.  Indeed, if
a formal edge index \(r\) satisfied \(1\le r<q-1\) and that exact edge
had color \(\beta\), then \noderef{TashkinovTree} would give an earlier
missing-\(\beta\) witness \(y_w\), with \(w\le r<q-1\).  This witness
and \(y_{q-1}\) are distinct by the displayed-vertex injectivity in
\noderef{TreeSequence}; the entry identities in
\noderef{TreeSequencePrefix} put both in the vertex set
\noderef{TreeSequenceVertexSet} of the proper prefix \(T_{q-1}\).
They both miss \(\beta\), contrary to the elementariness supplied by
\noderef{TashTreeProperPrefixElementary} and defined in
\noderef{ColoredVertexSetElementary}.  This proves precisely the
predicate \noderef{TreeColorUnusedThrough} for \(\beta\).

Apply \noderef{TashTreeControlledInterchange} to \(\alpha,\beta\) at
\(y_m,y_{q-1}\), and let \(\varphi'\) be the returned exact component
swap from \noderef{TwoColorComponentSwap}.  It makes \(\beta\) missing
at \(y_m\) and makes \((\varphi',T)\) a minimum bad pair.  Moreover,
\(y_n\) is outside the controlled component.  The returned component
is a path with the two distinct degree-at-most-one vertices
\(y_m,y_{q-1}\); they and \(y_n\) are pairwise distinct because their
indices are distinct and \noderef{TreeSequence} has an injective vertex
map.  The old missingness of \(\beta\) at \(y_n\), the properness from
\noderef{KTriple}, and \noderef{MissingColorIncidenceEndpoint} would
make \(y_n\) a third such vertex if it belonged to the component,
contradicting the spanning simple path supplied, with the same
properness premise, by \noderef{TwoColorComponentSpanningPath}.  The
outside clause of \noderef{TwoColorComponentSwapMissingBehavior}
therefore preserves missingness of \(\beta\) at \(y_n\), again using
the properness already extracted.  Since
\(m\le q-2\), this gives the required early index.

It remains to verify the no-repetition condition under \(\varphi'\).
Both ends \(y_m,y_{q-1}\) of the controlled path occur before \(q\).
For every suffix vertex \(y_s\), \(q\le s<n\), a vertex outside the
component keeps its complete missing set by
\noderef{TwoColorComponentSwapMissingBehavior}.  A suffix vertex inside
the component cannot be either end, so it is internal on the spanning
path from \noderef{TwoColorComponentSpanningPath}.  Such an internal
vertex has two distinct path neighbors.  By the bipartite adjacency
definition \noderef{TwoColorIncidenceGraph}, those neighbors are exact
incident edge-nodes colored \(\alpha\) or \(\beta\).  Properness from
\noderef{KTriple} prevents the two distinct incident edges from having
the same color, so there is one edge of each affected color.  The exact
transposition in \noderef{TwoColorComponentSwap} again leaves one of
each, and the other colors are unchanged.  Thus its complete missing set
is unchanged.  We already proved that the final vertex lies outside,
so its missing set is unchanged.  Colors other than
\(\alpha,\beta\) have unchanged missingness everywhere by the final
clause of \noderef{TwoColorComponentSwapMissingBehavior}.  Consequently
a color repeated between \(y_n\) and a suffix vertex after the swap
would already have been repeated there before the swap, contrary to the
hypothesis.  Hence \(\varphi',\beta,m\) have every asserted property.
\end{proof}
